
\documentclass{article}
\usepackage{colortbl}
\usepackage{makecell}
\usepackage{multirow}
\usepackage{supertabular}

\begin{document}

\newcounter{utterance}

\centering \large Interaction Transcript for game `matchit\_ascii', experiment `similar\_grid\_1', episode 4 with NEW\_Reflect\_Agent.
\vspace{24pt}

{ \footnotesize  \setcounter{utterance}{1}
\setlength{\tabcolsep}{0pt}
\begin{supertabular}{c@{$\;$}|p{.15\linewidth}@{}p{.15\linewidth}p{.15\linewidth}p{.15\linewidth}p{.15\linewidth}p{.15\linewidth}}
    \# & $\;$A & \multicolumn{4}{c}{Game Master} & $\;\:$B\\
    \hline

    \theutterance \stepcounter{utterance}  
    & & \multicolumn{4}{p{0.6\linewidth}}{
        \cellcolor[rgb]{0.9,0.9,0.9}{
            \makecell[{{p{\linewidth}}}]{
                \texttt{\tiny{[P1$\langle$GM]}}
                \texttt{You are participating in a collaborative guessing game.} \\
\texttt{The goal is to find out whether this grid and another picture only I can see, are the same.} \\
\texttt{Please describe your grid first. Then, I will provide my description and we can ask each other questions about the grids to figure out whether they are the same. Now start your short grid description with "DESCRIPTION:" followed by the description. Do not add anything else. This is your grid:} \\
\texttt{X ▢ ▢ ▢ X} \\
\texttt{▢ X ▢ X ▢} \\
\texttt{▢ ▢ X ▢ ▢} \\
\texttt{▢ ▢ X ▢ ▢} \\
\texttt{▢ ▢ X ▢ ▢} \\
            }
        }
    }
    & \\ \\

    \theutterance \stepcounter{utterance}  
    & \multicolumn{4}{p{0.45\linewidth}}{
        \cellcolor[rgb]{0.95,0.95,0.95}{
            \makecell[{{p{\linewidth}}}]{
                \texttt{\tiny{[P1$|$internal]}}
                \texttt{Reflections from previous episodes:} \\
\texttt{{-}} \\
\texttt{ Failed. I described the grid accurately but did not ask clarifying questions about the other grid to compare details. Next time, after giving my description, I should actively ask specific questions about the other grid's pattern to identify matches or mismatches.} \\
\texttt{{-}} \\
\texttt{ Failed. I gave a clear description of my grid pattern, but the user did not provide their description or engage in questioning, ending the episode prematurely. Next time, I would follow my initial description by immediately prompting for their description and proposing specific comparative questions to drive the collaboration forward.} \\
\texttt{{-}} \\
\texttt{ Failed. I described the grid as having a vertical line of X's in the third column, but I didn't specify the exact pattern of the first row (all X's). This could cause mismatch if the other grid differed in the top row. Next time, I would give a complete description including any full rows or columns, not just the central feature.} \\
            }
        }
    }
    & & \\ \\

    \theutterance \stepcounter{utterance}  
    & \multicolumn{4}{p{0.45\linewidth}}{
        \cellcolor[rgb]{0.95,0.95,0.95}{
            \makecell[{{p{\linewidth}}}]{
                \texttt{\tiny{[P1$|$internal]}}
                \texttt{{[}attempt 1 – valid{]}} \\
\texttt{DESCRIPTION: My grid has an X in the center of the third column, forming a vertical line of X's from the second row to the fifth row. The first row has X's in the first and fifth positions, with empty cells in between. The second row has X's in the second and fourth positions. The third, fourth, and fifth rows each have a single X in the third column, and all other cells are empty.} \\
            }
        }
    }
    & & \\ \\

    \theutterance \stepcounter{utterance}  
    & \multicolumn{4}{p{0.6\linewidth}}{
        \cellcolor[rgb]{0.9,0.9,0.9}{
            \makecell[{{p{\linewidth}}}]{
                \texttt{\tiny{[P1$\rangle$GM]}}
                \texttt{DESCRIPTION: My grid has an X in the center of the third column, forming a vertical line of X's from the second row to the fifth row. The first row has X's in the first and fifth positions, with empty cells in between. The second row has X's in the second and fourth positions. The third, fourth, and fifth rows each have a single X in the third column, and all other cells are empty.} \\
            }
        }
    }
    & & \\ \\

    \theutterance \stepcounter{utterance}  
    & & & \multicolumn{2}{p{0.3\linewidth}}{
        \cellcolor[rgb]{0.9,0.9,0.9}{
            \makecell[{{p{\linewidth}}}]{
                \texttt{\tiny{[GM$|$GM]}}
                \texttt{continue} \\
            }
        }
    }
    & & \\ \\

    \theutterance \stepcounter{utterance}  
    & & \multicolumn{4}{p{0.6\linewidth}}{
        \cellcolor[rgb]{0.9,0.9,0.9}{
            \makecell[{{p{\linewidth}}}]{
                \texttt{\tiny{[GM$\rangle$P2]}}
                \texttt{You are participating in a collaborative guessing game.} \\
\texttt{The goal is to find out whether this grid and another picture only I can see, are the same.} \\
\texttt{Please describe your grid first. Then, I will provide my description and we can ask each other questions about the grids to figure out whether they are the same. Now start your short grid description with "DESCRIPTION:" followed by the description. Do not add anything else. This is your grid:} \\
\texttt{▢ ▢ X ▢ ▢} \\
\texttt{▢ ▢ X ▢ ▢} \\
\texttt{▢ ▢ X ▢ ▢} \\
\texttt{▢ X ▢ X ▢} \\
\texttt{X ▢ ▢ ▢ X} \\
            }
        }
    }
    & \\ \\

    \theutterance \stepcounter{utterance}  
    & & & \multicolumn{4}{p{0.45\linewidth}}{
        \cellcolor[rgb]{0.95,0.95,0.95}{
            \makecell[{{p{\linewidth}}}]{
                \texttt{\tiny{[P2$|$internal]}}
                \texttt{Reflections from previous episodes:} \\
\texttt{{-}} \\
\texttt{ Failed. My description incorrectly placed the diamond's side Xs: I said second row fourth column, third row third, fourth row second, but my grid shows second row second column, third row third, fourth row fourth. This mismatch likely misled the user. Next time, I will double{-}check each coordinate against the actual grid before describing.} \\
\texttt{{-}} \\
\texttt{ FAIL. My description was inaccurate: I described a plus sign, but my grid has an extra X at (4,2) and (4,4), making the shape more like a cross with a thicker base or a small "T" shape, not a symmetric plus. This mismatch likely caused confusion. Next time, I would count X's precisely and describe exact coordinates of deviations from the simple plus, e.g., "plus sign with two additional X's symmetrically below the center row."} \\
\texttt{{-}} \\
\texttt{ Failed. My description incorrectly stated the first row was all X's, but my grid actually had only the third column as X in rows 1{-}4, with a full bottom row of X's. This major error misled the comparison from the start. Next time, I must describe the grid accurately and systematically, row by row, to avoid such a critical mistake.} \\
            }
        }
    }
     \\ \\

    \theutterance \stepcounter{utterance}  
    & & & \multicolumn{4}{p{0.45\linewidth}}{
        \cellcolor[rgb]{0.95,0.95,0.95}{
            \makecell[{{p{\linewidth}}}]{
                \texttt{\tiny{[P2$|$internal]}}
                \texttt{{[}attempt 1 – valid{]}} \\
\texttt{DESCRIPTION: A 5x5 grid with X's at positions (1,3), (2,3), (3,3), (4,2), (4,4), and (5,1), (5,5). All other cells are empty.} \\
            }
        }
    }
     \\ \\

    \theutterance \stepcounter{utterance}  
    & & & \multicolumn{4}{p{0.6\linewidth}}{
        \cellcolor[rgb]{0.9,0.9,0.9}{
            \makecell[{{p{\linewidth}}}]{
                \texttt{\tiny{[GM$\langle$P2]}}
                \texttt{DESCRIPTION: A 5x5 grid with X's at positions (1,3), (2,3), (3,3), (4,2), (4,4), and (5,1), (5,5). All other cells are empty.} \\
            }
        }
    }
     \\ \\

    \theutterance \stepcounter{utterance}  
    & & & \multicolumn{2}{p{0.3\linewidth}}{
        \cellcolor[rgb]{0.9,0.9,0.9}{
            \makecell[{{p{\linewidth}}}]{
                \texttt{\tiny{[GM$|$GM]}}
                \texttt{continue} \\
            }
        }
    }
    & & \\ \\

    \theutterance \stepcounter{utterance}  
    & & \multicolumn{4}{p{0.6\linewidth}}{
        \cellcolor[rgb]{0.9,0.9,0.9}{
            \makecell[{{p{\linewidth}}}]{
                \texttt{\tiny{[GM$\rangle$P2]}}
                \texttt{This is my DESCRIPTION: My grid has an X in the center of the third column, forming a vertical line of X's from the second row to the fifth row. The first row has X's in the first and fifth positions, with empty cells in between. The second row has X's in the second and fourth positions. The third, fourth, and fifth rows each have a single X in the third column, and all other cells are empty.} \\
\texttt{Now ask a question in order to find out new aspects of my grid that may be different to your grid. Start with "QUESTION:" and do not add anything else.} \\
            }
        }
    }
    & \\ \\

    \theutterance \stepcounter{utterance}  
    & & & \multicolumn{4}{p{0.6\linewidth}}{
        \cellcolor[rgb]{0.9,0.9,0.9}{
            \makecell[{{p{\linewidth}}}]{
                \texttt{\tiny{[GM$\langle$P2]}}
                \texttt{DESCRIPTION: X in the center of the third column, forming a vertical line of X's from the second row to the fifth row. The first row has X's in the first and fifth positions, with empty cells in between. The second row has X's in the second and fourth positions. The third, fourth, and fifth rows each have a single X in the third column, and all other cells are empty.  } \\
\texttt{QUESTION: Is the X in the first row located in the first and fifth columns?} \\
            }
        }
    }
     \\ \\

    \theutterance \stepcounter{utterance}  
    & & & \multicolumn{2}{p{0.3\linewidth}}{
        \cellcolor[rgb]{0.9,0.9,0.9}{
            \makecell[{{p{\linewidth}}}]{
                \texttt{\tiny{[GM$|$GM]}}
                \texttt{abort, first word: DESCRIPTION:, but expected QUESTION:} \\
            }
        }
    }
    & & \\ \\

\end{supertabular}
}

\end{document}
